\documentclass[11pt]{article}
\usepackage[margin=1in]{geometry}
\usepackage{amsmath,amssymb}
\usepackage{booktabs}
\usepackage{tabularx}
\usepackage{array}
\usepackage{enumitem}
\usepackage{hyperref}
\usepackage[T1]{fontenc}
\usepackage{caption}
\usepackage{xurl}

\hypersetup{colorlinks=true,linkcolor=blue,urlcolor=blue}
\setlength{\parindent}{0pt}
\setlength{\parskip}{0.55em}
\setlist[itemize]{leftmargin=*,itemsep=0.2em,topsep=0.3em}
\setlist[enumerate]{leftmargin=*,itemsep=0.2em,topsep=0.3em}
\captionsetup{font=small,skip=6pt}
\newcolumntype{Y}{>{\raggedright\arraybackslash}X}
\newcolumntype{P}[1]{>{\raggedright\arraybackslash}p{#1}}

\title{Rubric and weights (MUERP host score)\\[0.35em]
{\large companion to \texttt{muerp\_sequence\_findings.tex}}}
\author{Alexandra Bouillon}
\date{November 10, 2026}

\begin{document}
\maketitle

This is the scoring sheet used in \texttt{muerp\_sequence\_findings.tex}.
It rates a \textbf{simulator as a host for MUERP} (Zeng et~al., ICDCS 2024), not as a general quantum-network tool.

There are \textbf{15 metrics}, labeled \textbf{M1--M15}.
Each gets a score $s_i$ from 1 to 5, and a \textbf{weight} $w_i$ from its class.

\tableofcontents
\vspace{1em}

%----------------------------------------------------------------------
\section{The 1--5 scale}

Same meaning on every metric:

\begin{center}
\begin{tabularx}{\textwidth}{@{}r Y@{}}
\toprule
Score & Meaning \\
\midrule
5 & Native. Zeng et~al.'s construct maps 1:1 onto a documented hook. \\
4 & Works with modest extra code. No fight with the simulator. \\
3 & Possible, but you reimplement a piece Zeng et~al.\ treat as primitive. \\
2 & Workaround that breaks an assumption (for example shared qubits). \\
1 & Not practical without forking internals or changing the problem. \\
\bottomrule
\end{tabularx}
\end{center}

%----------------------------------------------------------------------
\section{Three classes, three weights}

Weights are not ``how hard the experiment was.''
They are ``how fatal a miss is if you want to host MUERP.''

\begin{center}
\begin{tabularx}{\textwidth}{@{}l l r Y@{}}
\toprule
Class & IDs & $w$ & What a miss means \\
\midrule
Must & M1--M6 & 3 & A 1 here means: do not use this simulator for MUERP. \\
Important & M7--M10 & 2 & Needed for an evaluation that looks like Zeng et~al., or for the timing model. \\
Nice & M11--M15 & 1 & Named extensions, inspectability, runtime, state representation. \\
\bottomrule
\end{tabularx}
\end{center}

Must-haves are the primitives in Zeng et~al.: users vs switches, exclusive $Q_r$, swap success $q$, length-dependent $p$, inject Algorithms~1--4, export slot $P$ (including $P=0$).

Nice-to-haves include two extensions Zeng et~al.\ \textbf{name} but do not plot: fidelity decay (M11) and concurrent groups sharing one $Q_r$ pool (M12).

\paragraph{Total weight.}
\begin{equation}
\sum_i w_i = 6\times 3 + 4\times 2 + 5\times 1 = 18 + 8 + 5 = 31.
\end{equation}

%----------------------------------------------------------------------
\section{Composite score}

If \textbf{any} Must-have is 1, stop. Do not average. The simulator is not a MUERP host.

Otherwise:
\begin{equation}
\mathrm{score} = \frac{\sum_i w_i s_i}{\sum_i w_i} = \frac{\sum_i w_i s_i}{31}.
\end{equation}

Always report \textbf{both}: the composite (one number), and the Must-have vector $[s_1,s_2,s_3,s_4,s_5,s_6]$.
A high average must not hide a Must-have 1.

\subsection{SeQUeNCe 1.3.1}

Must-have vector: $[3,5,5,4,4,3]$ --- no 1.

The findings line
\begin{verbatim}
Numerator 3(3 + 5 + 5 + 4 + 4 + 3) + 2(4 + 3 + 4 + 4) + (4 + 4 + 5 + 4 + 5) = 124.
Mean 124/31 = 4.
\end{verbatim}
is the composite written out.
Each parenthesis is one class of \textbf{scores} $s_i$.
The number in front is that class's \textbf{weight} $w$.

\begin{center}
\begin{tabularx}{\textwidth}{@{}l Y l@{}}
\toprule
Term & What it is & Arithmetic \\
\midrule
$3(3+5+5+4+4+3)$ & weight 3 $\times$ Must scores M1--M6 & $3\times 24=72$ \\
$2(4+3+4+4)$ & weight 2 $\times$ Important scores M7--M10 & $2\times 15=30$ \\
$(4+4+5+4+5)$ & weight 1 $\times$ Nice scores M11--M15 & $1\times 22=22$ \\
Numerator & $\sum_i w_i s_i$ & $72+30+22=124$ \\
31 & $\sum_i w_i=6\times 3+4\times 2+5\times 1$ & $18+8+5=31$ \\
Mean & numerator / 31 & $124/31=4$ \\
\bottomrule
\end{tabularx}
\end{center}

It is a \textbf{weighted} average, not $124/15$.
A Must-have point moves the mean three times as much as a nice-to-have point.

M10 is scored 4 as an interpolation: 3 would be ``timed only,'' 5 would be ``slotted \emph{and} timed.''
SeQUeNCe is timed; slot $P$ is Eq.~(2) after RSVP.

%----------------------------------------------------------------------
\section{Must-have (M1--M6), weight 3}

These six are what the findings section ``API map (M1--M6)'' walks through.
They are not the whole rubric.

{\small
\begin{center}
\begin{tabularx}{\textwidth}{@{}l P{1.55in} Y Y P{1.35in}@{}}
\toprule
ID & Question & 5 looks like & 3 looks like & SeQUeNCe \\
\midrule
M1 & Arbitrary graph; tag users $\neq$ switches? & Load any graph; first-class user vs switch types. & Custom graphs, but both roles are the same node type. & 3 --- both \texttt{QuantumRouter}. \\
M2 & Each channel through a switch reserves 2 exclusive qubits, and refuses when $Q_r$ is full? & Per-node memory, reservation, refusal. & Memories exist; exclusivity is your problem. & 5 --- Fig.~4a: $Q_S=4$ admits both, $Q_S=2$ rejects the second. \\
M3 & Two-qubit BSM swap with constant $q$? & Native swap; $q$ is a parameter. & Swap exists; $q$ is buried in a device model. & 5 --- \texttt{swapping\_success\_prob}. \\
M4 & Bell links with $p=\exp(-\alpha L)$? & Length and $\alpha$ are inputs. & Compute $p$ outside and paste it in. & 4 --- loss is dB/m; convert or compute $p$ after admission. \\
M5 & Algorithms~1--4 choose the tree; the simulator follows those paths? & Paths \emph{and} memories set before the run. & Paths can be set; the stack still picks memories. & 4 --- forwarding tables follow Algs~2--4; RSVP chooses memories. \\
M6 & Report Eq.~(1)--(2), including $P=0$ if a channel cannot be reserved? & Native $P$ from logs. & Reconstruct $P$ in a script. & 3 --- \texttt{muerp.eq2}; no SeQUeNCe metric named $P$. \\
\bottomrule
\end{tabularx}
\end{center}}

Failing Fig.~4a Case A or Case B fails M2 and/or M6, no matter how the rest looks.

%----------------------------------------------------------------------
\section{Important (M7--M10), weight 2}

\begin{center}
\begin{tabularx}{\textwidth}{@{}l Y Y@{}}
\toprule
ID & Question & SeQUeNCe \\
\midrule
M7 & Two channels may share fiber, not qubits. (That is why MUERP is not a Steiner tree.) & 4 --- two channels, 4 exclusive qubits on $S$. \\
M8 & Span all users with $|U|-1$ channels (a tree), not only isolated pairs. & 3 --- pairwise \texttt{RequestApp}s; the product $P$ is outside. \\
M9 & Run at Zeng et~al.'s size: 10 users + 50 switches (their eval also uses 20 graph replicas). & 4 --- one 10+50 replica, RSVP $\approx 0.5$\,s. \\
M10 & Either reproduce the simultaneous slot, or turn on wait/decoherence for the fidelity extension. & 4 --- timed; cutoff on/off; slot $P$ is Eq.~(2) after RSVP. \\
\bottomrule
\end{tabularx}
\end{center}

%----------------------------------------------------------------------
\section{Nice-to-have (M11--M15), weight 1}

\begin{center}
\begin{tabularx}{\textwidth}{@{}l Y Y@{}}
\toprule
ID & Question & SeQUeNCe \\
\midrule
M11 & Fidelity decay along a channel (named extension). & 4 --- swapping fidelity mean $0.57$ vs generation $0.90$ on the four-node layout. \\
M12 & Several MUERP groups sharing one $Q_r$ pool (named extension). & 4 --- two 5-user trees, both fit on that draw. \\
M13 & Inspect or change entanglement management, memory allocation, apps? & 5 --- open PyPI 1.3.1 modules. \\
M14 & Wall-clock and memory on \emph{this} workload (not a 1{,}000-node repeater prize). & 4 --- Waxman $0.55$\,s / $160$\,MB; chain $\le 2$\,s at 100 routers. \\
M15 & State representation cheap enough for many Bell pairs + BSM (not full wavefunctions). & 5 --- \texttt{bell\_diagonal} for this stack. \\
\bottomrule
\end{tabularx}
\end{center}

%----------------------------------------------------------------------
\section{How the pieces of the findings note relate}

\begin{center}
\begin{tabularx}{\textwidth}{@{}l Y@{}}
\toprule
In \texttt{muerp\_sequence\_findings.tex} & What it is \\
\midrule
Rating method & The scale, the three weights, the Fig.~4a gate. \\
API map (M1--M6) & Must-haves only: SeQUeNCe calls + those six scores. \\
Scorecard & All 15 metrics, weights, scores, evidence. \\
What these runs show & The composite $4.00$ and the Must-have vector. \\
\bottomrule
\end{tabularx}
\end{center}

The API map is not the rubric.
It is the Must-have slice, written as an interface note.
The rubric is M1--M15 plus the 1--5 anchors plus the weights $3/2/1$.

%----------------------------------------------------------------------
\section{What a 4.00 means}

SeQUeNCe 1.3.1 can host the MUERP \textbf{admission} model on the instances in the findings note.
It does not mean SeQUeNCe is ``more quantum,'' faster in general, or that Zeng et~al.'s algorithm ranking was re-checked.
A Must-have of 1 would have meant the opposite: do not host MUERP here.

\end{document}
